\chapter{Requirements, Design, and Architecture}

\section{Introduction}

This chapter translates the evidence and safety principles into an architecture
that can be implemented and tested. The design separates preparation of the
decoy from observation of an interaction, and it isolates format-dependent
behavior behind explicit policies.

\section{Actors and use cases}

Three actors interact with the core system:

\begin{itemize}
  \item the \textbf{defender/operator}, who configures JANUS, generates decoys,
  and investigates signals;
  \item the \textbf{local or remote actor}, who discovers, opens, modifies,
  moves, deletes, or uses a decoy;
  \item external \textbf{sensor services}, namely Wazuh and Canarytokens, which
  produce complementary evidence.
\end{itemize}

\plannedfigure{figures/system-context.png}
{JANUS system context and external actors}{fig:system-context}

\section{Functional requirements}

\begin{table}[H]
\centering
\caption{Main functional requirements}
\label{tab:functional-requirements}
\small
\begin{tabular}{p{0.12\linewidth}p{0.79\linewidth}}
\toprule
ID & Requirement \\
\midrule
FR-01 & Generate a HoneyDoc for an authorized business type and output format. \\
FR-02 & Select a compatible official token type and document its activation condition. \\
FR-03 & Reject unsafe, incoherent, or unsupported content before token creation and deployment. \\
FR-04 & Restrict deployment to the configured project-local root. \\
FR-05 & Register the artifact, token metadata, hash, TTL, and deployment path in SQLite. \\
FR-06 & Collect Wazuh alerts, correlate them with active HoneyDocs, and normalize the action. \\
FR-07 & Preserve raw evidence and a deterministic SHA-256 evidence hash. \\
FR-08 & Display local Wazuh signals and remote Canarytokens callbacks separately. \\
FR-09 & Support controlled token retirement and HoneyDoc rotation. \\
\bottomrule
\end{tabular}
\end{table}

\section{Non-functional requirements}

\begin{itemize}
  \item \textbf{Safety:} generated artifacts contain no real secret and no
  macro; paths remain within the deployment root.
  \item \textbf{Traceability:} every signal can be linked to a HoneyDoc and to
  immutable raw evidence.
  \item \textbf{Honesty:} fields distinguish direct observation, inference,
  confidence, and unavailable data.
  \item \textbf{Security:} administrative API routes require an API key;
  secrets are excluded from Git and routine logs.
  \item \textbf{Reproducibility:} a local content fallback and automated tests
  keep the project runnable without a remote LLM.
  \item \textbf{Maintainability:} generators, token policy, deployment,
  collection, correlation, and presentation remain separate modules.
\end{itemize}

\section{Global architecture}

The system is divided into five logical planes:

\begin{enumerate}[label=\arabic*.]
  \item \textbf{Control plane:} FastAPI endpoints, authentication, configuration,
  and the Streamlit operator interface.
  \item \textbf{Generation plane:} corpus-based prompt construction, LLM or
  fallback content, coherence checks, token policy, and artifact assemblers.
  \item \textbf{Deployment plane:} authorized-path resolution, atomic writes,
  hashing, TTL, and the SQLite registry.
  \item \textbf{Observation plane:} Windows SACL/event 4663, Wazuh rules and
  Indexer, remote Canarytokens, and optional Linux audit telemetry.
  \item \textbf{Evidence plane:} raw alerts, normalized events, verdicts,
  timelines, and dashboard views.
\end{enumerate}

\plannedfigure{figures/detection-architecture.png}
{End-to-end JANUS generation and detection architecture}
{fig:detection-architecture}

\section{Generation pipeline}

The API receives a document type, format, scenario, target subdirectory, and
TTL. The deployment target is resolved before prompt construction so an
untrusted path cannot escape the allowed root. The LLM receives only business
context. Content is checked before a remote token is allocated. A deterministic
assembler then creates the requested file and a final gate inspects its
structure before deployment.

\plannedfigure{figures/generation-pipeline.png}
{Controlled HoneyDoc generation pipeline}{fig:generation-pipeline}

\subsection{Supported format policy}

\begin{table}[H]
\centering
\caption{Document types, formats, and remote activation}
\label{tab:format-policy}
\small
\begin{tabular}{p{0.22\linewidth}p{0.25\linewidth}p{0.20\linewidth}p{0.24\linewidth}}
\toprule
Document type & Formats & Token & Remote trigger \\
\midrule
Financial report & DOCX, XLSX, CSV, JSON & Word, Excel, or web & Office external content or breadcrumb use \\
HR document & DOCX, CSV, JSON & Word or web & Office external content or breadcrumb use \\
Technical configuration & DOCX, ENV, YAML, JSON, ZIP & Word or web & Office external content or breadcrumb use \\
Cloud credentials & ENV, YAML, JSON, ZIP & AWS keys & Decoy key used with AWS API \\
\bottomrule
\end{tabular}
\end{table}

\section{Local detection design}

The monitored directory is \texttt{data/shared}. A SACL entry inherited by
children requests successful audit events for read, create/write, append, and
delete rights. The Wazuh agent collects the Windows Security channel and the
manager applies rules scoped to this path.

\begin{table}[H]
\centering
\caption{JANUS Wazuh rule taxonomy}
\label{tab:wazuh-rules}
\small
\begin{tabular}{p{0.15\linewidth}p{0.28\linewidth}p{0.48\linewidth}}
\toprule
Rule & Normalized meaning & Interpretation \\
\midrule
100100 & Generic audited access & A 4663 event targets the monitored decoy root. \\
100101 & Read/open candidate & ReadData or equivalent was exercised. \\
100102 & Modification & WriteData or AppendData was exercised. \\
100103 & Deletion & Delete rights were exercised. \\
100104 & Office/PDF open candidate & A relevant application accessed the decoy. \\
100110 & Windows logon context & Candidate session context for correlation. \\
100111 & SMB access context & Candidate remote file-share context. \\
100112 & Process context & Relevant Windows process execution. \\
100130 & Linux command context & auditd/execve command signal. \\
\bottomrule
\end{tabular}
\end{table}

\section{Interaction sequence}

The sequence begins when an actor touches an active decoy. Windows emits the
audit event, the agent forwards it, Wazuh classifies it, and the JANUS collector
stores both raw and normalized evidence. Independently, a supported remote
token may activate. These paths can have different delays and are not forced
into one synthetic event.

\plannedfigure{figures/interaction-sequence.png}
{Sequence of a HoneyDoc interaction and evidence collection}
{fig:interaction-sequence}

\section{Registry and evidence model}

SQLite stores four related concepts:

\begin{itemize}
  \item a \textbf{HoneyDoc} instance with filename, path, type, format, hash,
  creation time, TTL, and active state;
  \item a \textbf{token} with provider, type, activation mode, callback metadata,
  and lifecycle state;
  \item a \textbf{Wazuh detection} with the Wazuh document/event identifiers,
  rule identifier, normalized action, process, user, and verdict;
  \item \textbf{raw evidence and security signals} linked by a canonical
  SHA-256 hash and correlation keys.
\end{itemize}

\plannedfigure{figures/evidence-data-model.png}
{Simplified registry and evidence data model}{fig:evidence-data-model}

\section{Trust boundaries and safeguards}

The principal trust boundaries are the LLM API, public Canarytokens service,
Wazuh Indexer, monitored filesystem, and administrative API. The LLM never
receives token secrets or the deployment path. The public token service receives
only the information required to create a token. Indexer credentials and the
JANUS API key remain in \texttt{.env}. Administrative responses avoid exposing
auth tokens and cloud secret material. The public webhook remains disabled
until a real HTTPS endpoint can authenticate and receive it.

\section{Conclusion}

The architecture treats generation, deployment, and observation as separate
security domains connected by explicit data contracts. Chapter 4 describes the
implementation and shows that these contracts hold in automated and live tests.

